\documentclass{beamer}
%\documentclass[handout]{beamer}
%\documentclass[notes]{beamer}
%\documentclass[notes=only]{beamer}

%\ifdefined\handout
%\documentclass[handout,aspectratio=1610,xcolor={usenames,dvipsnames,table}]{beamer}
%\else
\documentclass[aspectratio=1610,xcolor={usenames,dvipsnames,table}]{beamer}
%\fi

\newcommand{\mylecture}{Configuration Management}

%\ifdefined\handout
%\documentclass[handout,aspectratio=1610,xcolor={usenames,dvipsnames,table}]{beamer}
%\else
\documentclass[aspectratio=1610,xcolor={usenames,dvipsnames,table}]{beamer}
%\fi

\newcommand{\mylecture}{Configuration Management}

%\ifdefined\handout
%\documentclass[handout,aspectratio=1610,xcolor={usenames,dvipsnames,table}]{beamer}
%\else
\documentclass[aspectratio=1610,xcolor={usenames,dvipsnames,table}]{beamer}
%\fi

\newcommand{\mylecture}{Configuration Management}

\input{../setup}
\input{../shared/setup}

\lstdefinelanguage{dump}
{
	morekeywords={kdbOpen,ksNew,keyNew,keyMeta,keyCopyMeta,keyEnd,ksEnd,kdbClose},
	sensitive=false,
	morecomment=[l]{//},
	morecomment=[s]{/*}{*/},
	morestring=[b]",
}


\lstdefinelanguage{SpecElektra}{
	%
	comment=[l]{;},
	commentstyle=\color{purple}\ttfamily,
	%
	morestring=[b]',
	morestring=[b]`,
	morestring=[b]",
	stringstyle=\color{purple}\ttfamily,
	%
	sensitive=f,% keywords are not case sensitive
	%
	% Colors see https://en.wikibooks.org/wiki/LaTeX/Colors
	%
	keywordstyle=\color{BlueViolet}\bfseries,
	keywordstyle=[2]\color{Green},
	keywordstyle=[3]\color{Aquamarine}\bfseries\textit,
	keywordstyle=[4]\color{NavyBlue}\bfseries,
	keywordstyle=[5]\color{Mahogany},
	%
	keywords={layer, require, validation, check, range, description, rationale, requirement, accessibility, enable, condition, message, default, opt, readonly, type, context, property1, property2, description, file, content, mountpoint, metadata, infos, plugins},
	keywords=[2]{},
	keywords=[3]{order, interface, network, emphasized},
	%keywords=[4]{[, ]},  %Not needed
	keywords=[4]{},
	keywords=[5]{},
	%
	literate={:=}{{{\color{red}\textbf:=}}}2
		 {\%}{{{\color{NavyBlue}\textbf\%}}}1
		 {[}{{{\color{Sepia}\textbf[}}}1
		 {]}{{{\color{Sepia}\textbf]}}}1,
}

\lstdefinelanguage{Cpp}{%
	language     = C++,
	literate=
}


\lstdefinelanguage{CfgElektra}{
	comment=[l]{;},
	commentstyle=\color{purple}\ttfamily,
	%
	morestring=[b]',
	morestring=[b]`,
	morestring=[b]",
	stringstyle=\color{purple}\ttfamily,
	%
	%
	sensitive=f,% keywords are not case sensitive
	%
	% Colors see https://en.wikibooks.org/wiki/LaTeX/Colors
	%
	keywordstyle=\color{Bittersweet}\bfseries,
	keywordstyle=[2]\color{DarkOrchid}\bfseries,
	keywordstyle=[3]\color{ForestGreen}\bfseries\textit,
	keywordstyle=[4]\color{Goldenrod}\bfseries,
	keywordstyle=[5]\color{CarnationPink},
	%
	keywords={},
	keywords=[2]{},
	keywords=[3]{},
	keywords=[4]{},
	keywords=[5]{},
	%
	literate={=}{{{\color{ForestGreen}\textbf=}}}1
		 %{<-}{{{\color{ForestGreen}\textbf<-}}}2
		 %{*}{{{\color{Bittersweet}\textbf*}}}1
		 {\%}{{{\color{NavyBlue}\textbf\%}}}1,
}




\lstset{language=SpecElektra, % Use SpecElektra as default programming language
	%boxpos=t, % make boxes a bit more unbreakable
	%frame=lines, % top+bottom line
	basicstyle=\ttfamily, % Use normal-size true type font
	showspaces,%
	showstringspaces=false, % Don't put marks in string spaces
	showlines=true, % make sure empty lines at end are shown (needed for concurrency
	tabsize=4, % spaces per tab
	xleftmargin=\parindent, % should be 18pt or 1.5em as defined by memoir
	%Does not really work well (needs to be deactivated for shortlistings):
	breaklines=false,
	%postbreak=\mbox{\textcolor{red}{$\hookrightarrow$}\space},
	%breakautoindent=true,
	%prebreak={\mbox{\ensuremath{\curvearrowright}}} % Zeichen am Zeilenende (Umbruch)
	%breaklines=true,
	%breakautoindent=true,
	%prebreak=\small\symbol{'134}, % backslash
	%prebreak={\mbox{\ensuremath{\curvearrowright}}} % lange kure
	%prebreak={\mbox{\ensuremath{\hookleftarrow}}} % lange kure
	%xleftmargin=3.0ex, %for some formats
	%xrightmargin=1.0ex, %for some formats
	%
	% Files do not work in utf8 see also:
	% http://stackoverflow.com/questions/1116266/listings-in-latex-with-utf-8-or-at-least-german-umlauts
	% http://tex.stackexchange.com/questions/24528/having-problems-with-listings-and-utf-8-can-it-be-fixed
	% Should work but doesn't? (Maybe add to literate broken?)
	%add to literate={ö}{{\"o}}1
	%	{ä}{{\"a}}1
	%	{ü}{{\"u}}1
	%	{Ö}{{\"O}}1
	%	{Ä}{{\"A}}1
	%	{Ü}{{\"U}}1
	%	{ß}{{\ss}}1,
	%
	% listingsutf8 did not work, made umlauts in comments very strange
	%extendedchars=true,
	%inputencoding=utf8,
	%
	%morecomment=[l][\color{blue}]{...}, % Line continuation (...) e.g. blue comment
	morekeywords={for_each},
	numbers=left, % Line numbers on left
	firstnumber=1, % Line numbers start with line 1
	numberstyle=\small\color{blue}, % Line numbers are blue and small
	numbersep=5pt,
	%stepnumber=5 % Line numbers go in steps of 5
}



\lstMakeShortInline[postbreak=,keywordstyle={}]^

\graphicspath{{../pic/}{../figures/}{../graphics/}{../ipe/}{../ggplot/}}



\lstdefinelanguage{dump}
{
	morekeywords={kdbOpen,ksNew,keyNew,keyMeta,keyCopyMeta,keyEnd,ksEnd,kdbClose},
	sensitive=false,
	morecomment=[l]{//},
	morecomment=[s]{/*}{*/},
	morestring=[b]",
}


\lstdefinelanguage{SpecElektra}{
	%
	comment=[l]{;},
	commentstyle=\color{purple}\ttfamily,
	%
	morestring=[b]',
	morestring=[b]`,
	morestring=[b]",
	stringstyle=\color{purple}\ttfamily,
	%
	sensitive=f,% keywords are not case sensitive
	%
	% Colors see https://en.wikibooks.org/wiki/LaTeX/Colors
	%
	keywordstyle=\color{BlueViolet}\bfseries,
	keywordstyle=[2]\color{Green},
	keywordstyle=[3]\color{Aquamarine}\bfseries\textit,
	keywordstyle=[4]\color{NavyBlue}\bfseries,
	keywordstyle=[5]\color{Mahogany},
	%
	keywords={layer, require, validation, check, range, description, rationale, requirement, accessibility, enable, condition, message, default, opt, readonly, type, context, property1, property2, description, file, content, mountpoint, metadata, infos, plugins},
	keywords=[2]{},
	keywords=[3]{order, interface, network, emphasized},
	%keywords=[4]{[, ]},  %Not needed
	keywords=[4]{},
	keywords=[5]{},
	%
	literate={:=}{{{\color{red}\textbf:=}}}2
		 {\%}{{{\color{NavyBlue}\textbf\%}}}1
		 {[}{{{\color{Sepia}\textbf[}}}1
		 {]}{{{\color{Sepia}\textbf]}}}1,
}

\lstdefinelanguage{Cpp}{%
	language     = C++,
	literate=
}


\lstdefinelanguage{CfgElektra}{
	comment=[l]{;},
	commentstyle=\color{purple}\ttfamily,
	%
	morestring=[b]',
	morestring=[b]`,
	morestring=[b]",
	stringstyle=\color{purple}\ttfamily,
	%
	%
	sensitive=f,% keywords are not case sensitive
	%
	% Colors see https://en.wikibooks.org/wiki/LaTeX/Colors
	%
	keywordstyle=\color{Bittersweet}\bfseries,
	keywordstyle=[2]\color{DarkOrchid}\bfseries,
	keywordstyle=[3]\color{ForestGreen}\bfseries\textit,
	keywordstyle=[4]\color{Goldenrod}\bfseries,
	keywordstyle=[5]\color{CarnationPink},
	%
	keywords={},
	keywords=[2]{},
	keywords=[3]{},
	keywords=[4]{},
	keywords=[5]{},
	%
	literate={=}{{{\color{ForestGreen}\textbf=}}}1
		 %{<-}{{{\color{ForestGreen}\textbf<-}}}2
		 %{*}{{{\color{Bittersweet}\textbf*}}}1
		 {\%}{{{\color{NavyBlue}\textbf\%}}}1,
}




\lstset{language=SpecElektra, % Use SpecElektra as default programming language
	%boxpos=t, % make boxes a bit more unbreakable
	%frame=lines, % top+bottom line
	basicstyle=\ttfamily, % Use normal-size true type font
	showspaces,%
	showstringspaces=false, % Don't put marks in string spaces
	showlines=true, % make sure empty lines at end are shown (needed for concurrency
	tabsize=4, % spaces per tab
	xleftmargin=\parindent, % should be 18pt or 1.5em as defined by memoir
	%Does not really work well (needs to be deactivated for shortlistings):
	breaklines=false,
	%postbreak=\mbox{\textcolor{red}{$\hookrightarrow$}\space},
	%breakautoindent=true,
	%prebreak={\mbox{\ensuremath{\curvearrowright}}} % Zeichen am Zeilenende (Umbruch)
	%breaklines=true,
	%breakautoindent=true,
	%prebreak=\small\symbol{'134}, % backslash
	%prebreak={\mbox{\ensuremath{\curvearrowright}}} % lange kure
	%prebreak={\mbox{\ensuremath{\hookleftarrow}}} % lange kure
	%xleftmargin=3.0ex, %for some formats
	%xrightmargin=1.0ex, %for some formats
	%
	% Files do not work in utf8 see also:
	% http://stackoverflow.com/questions/1116266/listings-in-latex-with-utf-8-or-at-least-german-umlauts
	% http://tex.stackexchange.com/questions/24528/having-problems-with-listings-and-utf-8-can-it-be-fixed
	% Should work but doesn't? (Maybe add to literate broken?)
	%add to literate={ö}{{\"o}}1
	%	{ä}{{\"a}}1
	%	{ü}{{\"u}}1
	%	{Ö}{{\"O}}1
	%	{Ä}{{\"A}}1
	%	{Ü}{{\"U}}1
	%	{ß}{{\ss}}1,
	%
	% listingsutf8 did not work, made umlauts in comments very strange
	%extendedchars=true,
	%inputencoding=utf8,
	%
	%morecomment=[l][\color{blue}]{...}, % Line continuation (...) e.g. blue comment
	morekeywords={for_each},
	numbers=left, % Line numbers on left
	firstnumber=1, % Line numbers start with line 1
	numberstyle=\small\color{blue}, % Line numbers are blue and small
	numbersep=5pt,
	%stepnumber=5 % Line numbers go in steps of 5
}



\lstMakeShortInline[postbreak=,keywordstyle={}]^

\graphicspath{{../pic/}{../figures/}{../graphics/}{../ipe/}{../ggplot/}}




\title[Configuration Management \hspace{25mm} \insertframenumber/\inserttotalframenumber]{Configuration Management}

%\subtitle{}
\institute{Institute of Information Systems Engineering, TU Wien}

\date{2.3.2018}

\begin{document}

\renewcommand{\enquote}[1]{\emph{``#1''}} % Cannot be done earlier

%%%%%%%%%%%%%%%%%%%%%%%%%%%%%%%
\begin{frame}
	\titlepage
	\doclicenseThis
\end{frame}

\begin{frame}
	\frametitle{Language}
	\begin{itemize}
		\item Slides are in English
		\item Papers are in English
		\item Book is in English
	\end{itemize}
\end{frame}

\begin{assignment}
	\frametitle{Language of the Talk?}
	\begin{task}
	Hands up if you prefer German.
	\end{task}
	Unanimous preference of German required, otherwise English.
\end{assignment}

%%%%%%%%%%%%%%%%%%%%%%%%%%%%%%%%%%%%%%%%%% 
\section{Motivation}
{
\usebackgroundtemplate{\includegraphics[width=\paperwidth]{pics/clouds.jpg}}%
\begin{frame}
	\frametitle{Misconfiguration}
	\begin{itemize}
		\item configuration is a user interface for system administrators
		\item \empha[misconfiguration]{misconfigurations}~\cite{yin2011empirical,su2007autobash,attariyan2010automating,xu2015systems}
			are a major cause of system failures~\cite{wool2004quantitative,oppenheimer2003internet,pertet2005causes}
		\item much time needed to fix misconfigurations~\cite{rabkin2011static,oppenheimer2003internet,yin2011empirical,mahajan2002bgp}
	\end{itemize}
\end{frame}
\begin{frame}
	\frametitle{No-Futz}
	\begin{itemize}
		\item \citet{holland2001nofutz}~defined \empha{futzing} to denote \enquote{tinkering or fiddling experimentally with something.}
		\item With \intro[no-futz computing]{no-futz computing} \citet{holland2001nofutz} mean \enquote{that futzing should be allowed, but should never be required.}
		\item currently configuration is error-prone and under-specified, \empha{futzing} is often required
	\end{itemize}
\end{frame}
}
\begin{assignment}
	\frametitle{First Assignment}
	\begin{itemize}
		\item Have you already experienced misconfiguration?
		\item Did you read about misconfiguration in the news?
	\end{itemize}
	\begin{task}
	Discuss with your neighbor and tell us the best story.
	\end{task}
\end{assignment}

%%%%%%%%%%%%%%%%%%%%%%%%%%%%%%%%%%%%%%%%%% 
\section{Content Overview}

\subsection{Terminology}
\begin{frame}
	\frametitle{Terminology}
\begin{definition}
\label{def:configuration-setting}
A \intro[configuration setting]{configuration setting},
or \intro[setting|see{configuration setting}]{setting} in short,
fulfills these properties:
\begin{enumerate}
\item
It is provided by the execution environment.
\item
It is \empha[consume]{consumed} by an application.
\item
It consists of a key, a configuration value, and potentially \empha{metadata}.
The \intro{configuration value}, or \intro[value|see{configuration value}]{value} in short, influences the application's behavior.
\item
It can be \empha[produce]{produced} by the maintainer, user, or system administrator of the software.
\end{enumerate}
\end{definition}
\end{frame}

%%%%%%%%%%%%%%%%%%%%%%%%%%%%%%%%%%%%%%%%%% 
\subsection{Requirements}
\begin{frame}
\begin{restatable}{requirement}{reqLegacy}
A configuration library must be able to integrate (legacy) systems and must fully support (legacy) configuration files.%
\label{req:legacy}
\end{restatable}

\begin{restatable}{requirement}{reqIntegration}
A configuration library must mitigate the configuration integration problem.%
\label{req:integration}
\end{restatable}

\begin{restatable}{requirement}{reqValidate}
Validation of configuration settings must happen systematically before the application is even started.
\end{restatable}
\end{frame}

\begin{assignment}
	\frametitle{Requirements}
	\begin{task}
	Discuss about requirements a configuration framework should fulfil.
	\end{task}
\end{assignment}

\subsection{Topics}
\begin{frame}
	Topic: \textit{reduction of misconfiguration and configuration duplicates} \\
	Misconfiguration obviously needs configuration.
	We will discuss the ideas of:
	\begin{itemize}
		\item when configuration is needed (decisions, user interface)
		\item how duplication can be avoided (generation of artifacts)
		\item configuration-less systems (auto-detection)
	\end{itemize}
\end{frame}

\begin{frame}
	Topic: \textit{sources of configuration}
	\begin{itemize}
		\item semi-structured data
		\item configuration file formats
		\item command-line arguments
		\item environment variables
	\end{itemize}
\end{frame}

\begin{frame}
	Topic: \textit{desirable properties of configuration}
	\begin{itemize}
		\item self-description
		\item changeability
		\item idempotence
		\item round-tripping
	\end{itemize}
\end{frame}

\begin{frame}
	Topic: \textit{configuration as user interface}
	\begin{itemize}
		\item How system administrators work.
		\item Which user interfaces exist.
		\item Which structures exist.
	\end{itemize}
\end{frame}

\begin{frame}
	Topic: \textit{design of configuration and configuration access}
	\begin{itemize}
		\item architecture
		\item introspection
		\item code generation
	\end{itemize}
\end{frame}

\begin{assignment}
	\begin{task}
	Break.
	\end{task}
\end{assignment}

\begin{frame}
	Topic: \textit{influence of variability on software architecture}
	\begin{itemize}
		\item variability
		\item architectural decisions
	\end{itemize}
\end{frame}

\begin{frame}
	Topic: \textit{complexity reduction and testability of configurable applications}
	Are all configuration settings:
	\begin{itemize}
		\item used?
		\item tested?
		\item ...?
	\end{itemize}
\end{frame}

\begin{frame}
	Topic: \textit{early detection of misconfiguration}
	\begin{itemize}
		\item points in time for configuration access and validation
		\item validation techniques
	\end{itemize}
\end{frame}

\begin{frame}
	Topic: \textit{configuration management tools}
	\begin{itemize}
		\item Puppet
		\item CfEngine
		\item Nix
		\item others...
	\end{itemize}
\end{frame}

\begin{frame}
	Further topics:
	\begin{itemize}
		\item \textit{context-awareness, avoidance of dependences}
		\item \textit{cascading configuration}
		\item \textit{strategies for validation and modularization}
		\item \textit{documentation of configuration}
	\end{itemize}
\end{frame}

\begin{assignment}
	\frametitle{In which topics are you interested?}
	\begin{task}[1]
	Choose a partner for this task.
	\end{task}

	\begin{task}[2]
	Go to stations and discuss topic with your partner.
	\end{task}

	\begin{task}[3]
	Write down the most interesting topics.
	(Can be topics of stations or new topics.)
	\end{task}
\end{assignment}

\begin{frame}
	\frametitle{Elektra}
	\hfill \includegraphics[width=2cm]{../figures/logo}
	\vspace{-1cm}
	%Elektra is an implementation of the next generation of configuration systems:
	\begin{itemize}
		\item Elektra improves configuration management.
		\item Configuration management tools can use Elektra.
		\item Elektra implements what we discuss in this lecture.
		\item Elektra allows applications to fulfil the requirements we will discuss.
		\item Obviously you will also be able to apply the knowledge from this LVA without Elektra.
	\end{itemize}
\end{frame}





%%%%%%%%%%%%%%%%%%%%%%%%%%%%%%%%%%%%%%%%%% 
\section{Organisation}

\begin{frame}
	Communication:
	\begin{itemize}
		\item TISS (forum and news) \url{https://tiss.tuwien.ac.at/course/courseDetails.xhtml?courseNr=194030&semester=2018S}
		\item GitHub (private and public repo) \url{https://git.libelektra.org}
		\item EMail \url{markus.raab@complang.tuwien.ac.at}
		\item before/after/during lectures
	\end{itemize}
\end{frame}

\begin{assignment}
	\begin{task}
	Send me your GitHub name per email to get access to private repo.
	\end{task}
\end{assignment}

\subsection{Previous Knowledge}
\begin{frame}
	\begin{itemize}
		\item Obviously \textit{no} prior knowledge about Configuration or Configuration Management necessary.
		\item If you already have experience, you can use it in your talk and assignments.
		\item Knowledge about software engineering and software requirements is beneficially.
		\item You should have an understanding of large-scale software construction.
		\item Programming skills is a must.
	\end{itemize}
\end{frame}

\subsection{Grades}
\begin{frame}
	You will get a grade only if:
	\begin{itemize}
		\item You submitted your homework. (PR is not enough.)
		\item You participated in the team exercise.
		\item You gave your talk.
	\end{itemize}
\end{frame}

\subsection{Grades}
\begin{frame}
	To get a positive grade:
	\begin{itemize}
		\item All parts must be done.
		\item All parts must be positive.
	\end{itemize}
\end{frame}

\begin{frame}
	Grade is calculated from:
	\begin{description}
	\item[30\,\%:] homework
	\item[30\,\%:] team exercise
	\item[10\,\%:] talk
	\item[30\,\%:] test
	\item[+:] extrapoints
	\end{description}
\end{frame}

\subsection{Assignments}
\begin{frame}
	\frametitle{Talk}
	You can give a talk about anything related to configuration management.
	\begin{itemize}
		\item 20 minutes.
		\item It must be about your experience.
		\item I.e., not only about study of literature.
		\item It is okay if the experience happened during this LVA.
		\item If you extensively use some tool, please share your knowledge.
	\end{itemize}
\end{frame}

\begin{frame}
	\frametitle{Home Work}
	You can select your own task related to \emph{use} configuration management.
	For example:
	\begin{itemize}
		\item Use Elektra in a small application.
		\item Write a tutorial about how you did it.
		\item Fix usage bugs in Elektra, make it easier. (Extrapoints)
	\end{itemize}
\end{frame}

\begin{assignment}
	\begin{task}
	Think about a talk and task till next week.
	Write it down in the private repo (first come, first served).
	\end{task}
\end{assignment}

\begin{frame}
	\frametitle{Team Exercise}
	You can select your team (2-3 people) and your task related to \emph{improve} configuration management.
	For example, you can improve Elektra:
	\begin{itemize}
		\item Write a validation plugin.
		\item Write support for a configuration file format.
		\item Make a configuration management tool use Elektra.
		\item Fix bugs. (Extrapoints)
	\end{itemize}
\end{frame}

\begin{assignment}
	\begin{task}
	Talk with someone who is not your neighbor about a potential collaboration in the team exercise.
	\end{task}
\end{assignment}

\begin{assignment}
	\frametitle{Questions?}
	\begin{task}
	Please read TISS and register for the course.
	\end{task}

	\begin{task}
	Any questions?
	\end{task}
\end{assignment}



%%%%%%%%%%%%%%%%%%%%%%%%%%%%%%%%%%%%%%%%%% 
\subsection{Literature}

\nocite{raab2017introducing}

\begin{frame}[allowframebreaks]
	\bibliographystyle{plainnat}
	\bibliography{../shared/elektra.bib}
\end{frame}

\end{document}


